\chapter{Conception}
\label{chap:conception}

\section{Architecture globale}

L'application est découpée en six modules C, plus une bibliothèque
\emph{vendor} pour le parsing TOML. Chaque module possède un fichier
d'en-tête (dans \file{include/}) qui expose son API publique, et un
fichier source (dans \file{src/}) qui contient l'implémentation.
La figure~\ref{fig:architecture} présente le graphe des inclusions.

\begin{figure}[h!]
\centering
% Architecture des modules : graphe de dépendance d'inclusion.
\begin{tikzpicture}[
  node distance = 1.4cm and 1.6cm,
  module/.style = {
    draw, rounded corners, thick,
    minimum width = 2.2cm, minimum height = 0.9cm,
    align = center, font = \small\ttfamily,
    fill = blue!8
  },
  vendor/.style = {
    draw, rounded corners, thick, dashed,
    minimum width = 2.2cm, minimum height = 0.9cm,
    align = center, font = \small\ttfamily,
    fill = orange!10
  },
  arr/.style = {-Stealth, thick, gray!70}
]

% Niveau 1 : main
\node[module] (main) {main};

% Niveau 2 : world, conf
\node[module, below left = of main]  (world) {world};
\node[module, below right = of main] (conf)  {conf};

% Niveau 3 : entity, species
\node[module, below = of world] (entity)  {entity};
\node[module, below = of conf]  (species) {species};

% Niveau 4 : vec2, tomlc99
\node[module, below = of entity] (vec2) {vec2};
\node[vendor, below = of species] (toml) {tomlc99\\(vendor)};

% Inclusions
\draw[arr] (main)    -- (world);
\draw[arr] (main)    -- (conf);
\draw[arr] (world)   -- (entity);
\draw[arr] (world)   -- (species);
\draw[arr] (entity)  -- (species);
\draw[arr] (entity)  -- (vec2);
\draw[arr] (world)   -- (vec2);
\draw[arr] (conf)    -- (species);
\draw[arr] (conf)    -- (toml);

% Légende
\node[draw=none, font=\footnotesize, below = 1cm of vec2,
      anchor=north west, xshift=-0.5cm] (legend) {
  \begin{tabular}{l}
    \tikz \draw[-Stealth, thick, gray!70] (0,0) -- (0.5,0); \ \emph{inclusion}\\[2pt]
  \end{tabular}
};

\end{tikzpicture}

\caption{Architecture modulaire et dépendances d'inclusion.}
\label{fig:architecture}
\end{figure}

Rôle de chaque module~:

\begin{description}[leftmargin=3em, style=nextline]
  \item[\code{main}.] Point d'entrée de l'application. Crée la
    fenêtre GTK, met en place les contrôles latéraux, charge la
    configuration via \code{conf}, initialise le \code{World} et
    enregistre le \emph{tick} via \func{g\_timeout\_add}.

  \item[\code{world}.] Possède le pool d'entités, gère leur
    spawn/despawn, orchestre la séquence de mise à jour
    (déplacement, alimentation, vieillissement, respawn), et
    expose la fonction \func{world\_tick} appelée à chaque
    image.

  \item[\code{entity}.] Représente une entité individuelle.
    Implémente les forces de pilotage (boids, fuite, chasse,
    évitement intra-espèce), l'intégration physique
    semi-implicite, le lissage de l'angle de cap, et la mise à
    jour visuelle (positionnement et CSS).

  \item[\code{species}.] Registre dynamique des
    \struct{SpeciesConfig}. Stocke en mémoire le tableau des
    espèces issu de la configuration et fournit les fonctions
    de recherche par index ou par identifiant.

  \item[\code{conf}.] Lit \file{aquarium.conf} via le parseur
    TOML, valide le contenu, résout les listes \code{prey} en
    bitmask, et publie les paramètres globaux dans des variables
    \code{extern}.

  \item[\code{vec2}.] Bibliothèque \emph{header-only} de
    vecteurs 2D~: addition, soustraction, multiplication par
    scalaire, normalisation, limitation, \emph{wrapping} torique
    pour le monde, et calcul de delta minimal en distance
    enveloppée.

  \item[\code{tomlc99} (vendor).] Parseur TOML en C99 distribué
    en deux fichiers (\file{toml.h}, \file{toml.c}), embarqué
    dans \file{vendor/tomlc99/}.
\end{description}

\section{Modèle de données}

Trois structures principales portent l'état du système.

\subsection{\struct{Entity}}

Une \struct{Entity} représente une instance vivante (ou en
sursis) dans l'aquarium. Elle agrège trois groupes de champs~:

\begin{itemize}
  \item \textbf{Identité et appartenance}~: \code{id} unique,
    \code{species} (index dans le registre).
  \item \textbf{État physique}~: position \code{pos}, vitesse
    \code{vel}, accélération \code{acc}, angle de cap
    \code{angle} (lissé), facteur de variation individuel
    \code{speed\_scale}.
  \item \textbf{Cycle de vie et rendu}~: \code{state} (\code{ALIVE},
    \code{DYING}, \code{INACTIVE}), \code{age},
    \code{last\_eaten}, \code{respawn\_at}, \code{death\_timer},
    plus le \code{widget} GTK et le \code{css\_provider}
    associés.
\end{itemize}

\subsection{\struct{SpeciesConfig}}

Cette structure regroupe l'intégralité des paramètres
gouvernant le comportement d'une espèce~: identifiant
\code{id}, nom affiché, chemin du sprite, dimensions du sprite,
ratio de bouche, vitesses min/max, force max, rayons et poids
des trois forces de Reynolds, rayon et poids de fuite, rayon
et poids de chasse, rayon et poids d'évitement intra-espèce,
masque de proie, durée de vie, taille initiale du banc,
plafond de population, délai de respawn, et drapeau
\code{is\_apex}.

L'ensemble des paramètres comportementaux est ainsi capturé
dans un \emph{seul} type, ce qui simplifie considérablement
le code consommateur~: \func{entity\_tick} reçoit un pointeur
\code{const SpeciesConfig *} et lit les paramètres dont il a
besoin, sans avoir à connaître le système de configuration.

\subsection{\struct{World}}

Le \struct{World} contient le pool d'entités en tableau
statique \code{Entity entities[MAX\_ENTITIES]}, le compteur
courant, les dimensions de la zone simulée, le pointeur vers
le conteneur GTK, le pas de temps courant \code{dt}, le temps
écoulé \code{elapsed}, le générateur d'identifiants
\code{next\_id}, ainsi que deux tableaux indexés par espèce~:
\code{alive\_counts[MAX\_SPECIES]} et
\code{respawn\_ready[MAX\_SPECIES]}.

\section{Séquence d'un tick}

À chaque pas de simulation, \func{world\_tick} orchestre une
séquence ordonnée d'étapes. La figure~\ref{fig:tick-sequence}
représente cet enchaînement.

\begin{figure}[h!]
\centering
\resizebox{0.85\textwidth}{!}{% Diagramme de séquence d'un tick.
\begin{tikzpicture}[
  every node/.style = {font=\small},
  lifeline/.style   = {dashed, gray!60},
  actor/.style      = {draw, fill=blue!10, rounded corners,
                       minimum width=2.2cm, minimum height=0.7cm,
                       align=center, font=\small\ttfamily},
  msg/.style        = {-Stealth, thick},
  ret/.style        = {-Stealth, thick, dashed},
  activation/.style = {draw, fill=blue!20, minimum width=0.25cm}
]

% Acteurs (en haut)
\node[actor] (main)    at (0,0)  {main loop};
\node[actor] (world)   at (4,0)  {World};
\node[actor] (entity)  at (8,0)  {Entity[i]};
\node[actor] (species) at (12,0) {Species};

% Lignes de vie
\foreach \n in {main, world, entity, species} {
  \draw[lifeline] (\n.south) -- ++(0,-13);
}

% Messages (de haut en bas)
\draw[msg] ($(main.south)+(0,-0.5)$) -- node[above, font=\footnotesize]{tick\_cb}
           ($(world.south)+(0,-0.5)$);

\draw[msg] ($(world.south)+(0,-1.0)$) -- node[above, font=\footnotesize]{world\_tick}
           ($(world.south)+(2.5,-1.0)$) -- ++(0,-0.4) -| ($(world.south)+(0,-1.5)$);

% Boucle entity_tick
\draw[msg] ($(world.south)+(0,-2.2)$) -- node[above, font=\footnotesize]{entity\_tick}
           ($(entity.south)+(0,-2.2)$);
\draw[ret] ($(entity.south)+(0,-2.5)$) -- ($(world.south)+(0,-2.5)$);
\node[font=\footnotesize, gray] at (8, -3.0) {(boucle sur N entités)};

% Eating
\draw[msg] ($(world.south)+(0,-3.7)$) -- node[above, font=\footnotesize]{world\_handle\_eating}
           ($(world.south)+(2.0,-3.7)$) -- ++(0,-0.4) -| ($(world.south)+(0,-4.2)$);

% Lifespan
\draw[msg] ($(world.south)+(0,-4.7)$) -- node[above, font=\footnotesize]{world\_handle\_lifespans}
           ($(world.south)+(2.5,-4.7)$) -- ++(0,-0.4) -| ($(world.south)+(0,-5.2)$);

% Visuals
\draw[msg] ($(world.south)+(0,-6.0)$) -- node[above, font=\footnotesize]{entity\_apply\_visuals}
           ($(entity.south)+(0,-6.0)$);
\draw[ret] ($(entity.south)+(0,-6.3)$) -- ($(world.south)+(0,-6.3)$);
\node[font=\footnotesize, gray] at (8, -6.8) {(boucle sur N entités)};

% Respawn
\draw[msg] ($(world.south)+(0,-7.6)$) -- node[above, font=\footnotesize]{world\_handle\_respawns}
           ($(world.south)+(2.5,-7.6)$) -- ++(0,-0.4) -| ($(world.south)+(0,-8.1)$);

% Lookup species
\draw[msg] ($(world.south)+(0,-8.8)$) -- node[above, font=\footnotesize]{species\_get(s)}
           ($(species.south)+(0,-8.8)$);
\draw[ret] ($(species.south)+(0,-9.1)$) -- ($(world.south)+(0,-9.1)$);

% Retour final
\draw[ret] ($(world.south)+(0,-10.0)$) -- node[above, font=\footnotesize]{G\_SOURCE\_CONTINUE}
           ($(main.south)+(0,-10.0)$);

\end{tikzpicture}
}
\caption{Séquence d'un tick de simulation.}
\label{fig:tick-sequence}
\end{figure}

Chronologiquement~:

\begin{enumerate}
  \item Avancement du temps~: \code{elapsed += dt}.
  \item Pour chaque entité \code{ALIVE}, calcul des forces et
    mise à jour de la position et de la vitesse via
    \func{entity\_tick}.
  \item Détection et application des \emph{repas}~:
    \func{world\_handle\_eating} parcourt tous les couples
    prédateur/proie compatibles et déclenche un
    \func{world\_despawn} si la bouche du prédateur entre dans
    le rayon de hit de la proie.
  \item Vérification de la famine~: une entité qui n'a pas
    mangé depuis \code{STARVATION\_SECONDS} secondes est
    \emph{despawn}ée.
  \item Pour chaque entité, \func{entity\_apply\_visuals}
    actualise la position du widget et le CSS de rotation. Si
    l'entité est en \code{DYING}, le timer décroît jusqu'à
    \code{INACTIVE}.
  \item Tentative de respawn pour chaque espèce dont
    \code{alive\_counts} est inférieur à \code{flock\_size}.
\end{enumerate}

L'ordre des étapes est important. La détection des repas
\emph{après} l'intégration physique garantit que la collision
est calculée sur les positions à jour. La mise à jour visuelle
\emph{après} la détection des repas permet d'afficher le flash
de mort dès le tick suivant la consommation. Le respawn en
dernière étape évite de réactiver une entité qui aurait été
mangée le même tick.

\section{Cycle de vie d'une entité}

La figure~\ref{fig:lifecycle} formalise les trois états et
leurs transitions.

\begin{figure}[h!]
\centering
% Cycle de vie d'une entité : automate à 3 états.
\begin{tikzpicture}[
  ->, >=Stealth, thick,
  node distance = 4.5cm,
  every state/.style = {
    minimum size = 1.8cm,
    align = center,
    font = \small\ttfamily,
    draw = black,
    thick
  },
  alive/.style    = {state, fill=green!15},
  dying/.style    = {state, fill=orange!25},
  inactive/.style = {state, fill=gray!20}
]

\node[alive]    (alive)    {ALIVE};
\node[dying, right of=alive]    (dying)    {DYING};
\node[inactive, right of=dying] (inactive) {INACTIVE};

% Self-loop sur ALIVE (tick normal)
\path (alive) edge[loop above]
  node[font=\footnotesize, align=center]{tick normal\\(boids, manger, fuir)} (alive);

% ALIVE -> DYING : world_despawn
\path (alive) edge[bend left=15]
  node[above, font=\footnotesize, align=center]{world\_despawn() \\ (mangé / famine)} (dying);

% DYING -> INACTIVE : death_timer expiré
\path (dying) edge[bend left=15]
  node[above, font=\footnotesize, align=center]{death\_timer\\$\leq 0$} (inactive);

% INACTIVE -> ALIVE : respawn
\path (inactive) edge[bend left=45]
  node[below, font=\footnotesize, align=center]{respawn\_ready $\leq$ elapsed\\$\land$ alive\_count $<$ flock\_size} (alive);

\end{tikzpicture}

\caption{Machine à états du cycle de vie d'une entité.}
\label{fig:lifecycle}
\end{figure}

\begin{description}[leftmargin=3em, style=nextline]
  \item[\code{ALIVE} $\to$ \code{ALIVE}.] Cas nominal. L'entité
    est mise à jour à chaque tick.

  \item[\code{ALIVE} $\to$ \code{DYING}.] Déclenché par
    \func{world\_despawn}, qui peut être appelé soit par
    \func{world\_handle\_eating} (mort par prédation) soit par
    \func{world\_handle\_lifespans} (mort par famine). Le
    compteur \code{alive\_counts[species]} est décrémenté, et
    \code{death\_timer} est initialisé à
    \code{DEATH\_FLASH\_DURATION}.

  \item[\code{DYING} $\to$ \code{INACTIVE}.] À chaque tick en
    \code{DYING}, \code{death\_timer} décroît de \code{dt}.
    Lorsqu'il atteint zéro, le widget est masqué et l'état
    passe à \code{INACTIVE}.

  \item[\code{INACTIVE} $\to$ \code{ALIVE}.] Si la population
    vivante est inférieure à \code{flock\_size} \emph{et} que
    \code{elapsed} dépasse \code{respawn\_ready[species]},
    \func{world\_handle\_respawns} réactive une entité
    \code{INACTIVE} de cette espèce~: position et vitesse
    aléatoires, état \code{ALIVE}, widget rendu visible, et
    \code{respawn\_ready} repoussé de \code{respawn\_delay}.
\end{description}

Le choix d'un \emph{recyclage} d'entités plutôt qu'une
allocation/destruction permet de borner statiquement le tableau
\code{entities} à \code{MAX\_ENTITIES} et d'éviter toute
fragmentation~: les widgets GTK sont créés une fois pour toutes
au moment du \emph{spawn}, puis réutilisés.

\section{Graphe prédateur/proie}

La configuration par défaut implémente une chaîne alimentaire à
trois niveaux représentée dans la figure~\ref{fig:predator-prey}.

\begin{figure}[h!]
\centering
% Graphe prédateur/proie : les flèches vont du prédateur vers sa proie.
\begin{tikzpicture}[
  node distance = 1.4cm and 1.0cm,
  fish/.style = {
    draw, rounded corners, thick,
    minimum width = 1.6cm, minimum height = 0.7cm,
    align = center, font = \small
  },
  prey/.style    = {fish, fill = green!12},
  midpred/.style = {fish, fill = yellow!18},
  apex/.style    = {fish, fill = red!18, very thick},
  hunt/.style    = {-Stealth, thick}
]

% Apex
\node[apex] (shark) at (0, 4) {Shark};

% Prédateurs intermédiaires
\node[midpred] (barra) at (-4.5, 1.5) {Barracuda};
\node[midpred] (bass)  at (-1.5, 1.5) {Bass};
\node[midpred] (trout) at ( 1.5, 1.5) {Trout};
\node[midpred] (tuna)  at ( 4.5, 1.5) {Tuna};

% Proies
\node[prey] (gold)  at (-3, -1.2) {Goldfish};
\node[prey] (clown) at ( 0, -1.2) {Clownfish};
\node[prey] (angel) at ( 3, -1.2) {Angelfish};

% Shark mange Tuna, Bass, Trout, Angelfish, Goldfish
\draw[hunt, red!60] (shark) -- (tuna);
\draw[hunt, red!60] (shark) -- (bass);
\draw[hunt, red!60] (shark) -- (trout);
\draw[hunt, red!60] (shark) -- (angel);
\draw[hunt, red!60] (shark) to[bend right=20] (gold);

% Barracuda et Bass mangent Goldfish + Clownfish
\draw[hunt] (barra) -- (gold);
\draw[hunt] (barra) -- (clown);
\draw[hunt] (bass)  -- (gold);
\draw[hunt] (bass)  -- (clown);

% Trout mange Goldfish
\draw[hunt] (trout) -- (gold);

% Tuna mange Goldfish + Clownfish
\draw[hunt] (tuna) -- (gold);
\draw[hunt] (tuna) -- (clown);

% Légende
\matrix [draw, below right=0.6cm and -1.5cm of angel.south,
         column sep=4pt, row sep=2pt, font=\footnotesize] {
  \node[apex,    minimum width=0.8cm, minimum height=0.4cm] {}; & \node{Apex}; \\
  \node[midpred, minimum width=0.8cm, minimum height=0.4cm] {}; & \node{Prédateur intermédiaire}; \\
  \node[prey,    minimum width=0.8cm, minimum height=0.4cm] {}; & \node{Proie}; \\
};

\end{tikzpicture}

\caption{Chaîne alimentaire de l'écosystème par défaut.}
\label{fig:predator-prey}
\end{figure}

Le requin (\code{shark}) est l'unique prédateur \emph{apex}
(drapeau \code{is\_apex = true}, aucun \code{fear\_radius}). Il
chasse cinq espèces. Quatre prédateurs intermédiaires
(\code{barracuda}, \code{bass}, \code{trout}, \code{tuna})
chassent les petites proies (\code{goldfish}, \code{clownfish})
et craignent le requin. Trois espèces de proies pures
(\code{goldfish}, \code{clownfish}, \code{angelfish}) ne
chassent personne.

L'angelfish a un statut particulier~: il n'est chassé que par
le requin, ce qui en fait une proie résiliente, intermédiaire
entre les petites proies (massivement chassées) et les
prédateurs.

\section{Forces de pilotage}

La figure~\ref{fig:boids-forces} illustre comment les trois
forces de Reynolds s'appliquent à un agent. Les forces
supplémentaires (fuite, chasse, évitement intra-espèce) sont
construites sur le même schéma vectoriel mais avec des
voisinages et des cibles différents.

\begin{figure}[h!]
\centering
% Diagramme des trois forces de Reynolds (separation, alignment, cohesion)
% appliquées à un poisson central.
\begin{tikzpicture}[
  >=Stealth,
  fish/.style    = {circle, draw, fill=blue!15, minimum size=0.55cm, inner sep=0pt},
  selfish/.style = {circle, draw, fill=blue!40, minimum size=0.7cm, inner sep=0pt, very thick},
  sep/.style     = {-Stealth, very thick, red!80!black},
  ali/.style     = {-Stealth, very thick, green!50!black},
  coh/.style     = {-Stealth, very thick, blue!70!black}
]

% Poisson central (self)
\node[selfish] (s) at (0,0) {};

% Voisins
\node[fish] (n1) at (-1.4, 0.3) {};
\node[fish] (n2) at (-0.4, 1.3) {};
\node[fish] (n3) at ( 1.0, 1.0) {};
\node[fish] (n4) at ( 1.5,-0.4) {};
\node[fish] (n5) at ( 0.4,-1.3) {};
\node[fish] (n6) at (-1.0,-1.1) {};

% Vitesses des voisins (pour l'alignement)
\foreach \n in {n1, n2, n3, n4, n5, n6} {
  \draw[->, gray!50, thin] (\n.center) -- ++(0.6, 0.2);
}

% Vecteur de séparation (s'éloigner des voisins proches)
\draw[sep] (s.center) -- ++(-1.2, -0.6) node[below left, font=\small, red!80!black] {séparation};

% Vecteur d'alignement (suivre la moyenne des vitesses)
\draw[ali] (s.center) -- ++(1.4, 0.6) node[above right, font=\small, green!50!black] {alignement};

% Vecteur de cohésion (vers le centre du groupe)
\draw[coh] (s.center) -- ++(0.3, 1.3) node[above, font=\small, blue!70!black] {cohésion};

% Cercle de voisinage
\draw[dashed, gray!50] (s) circle (2.0);
\node[font=\footnotesize, gray] at (1.6, -1.7) {rayon de voisinage};

% Légende des formules en bas
\node[draw, font=\footnotesize, align=left,
      below = 0.6cm of s, anchor=north] (legend) {
  $\vec{F}_{\text{sep}} \propto -\sum_i \dfrac{\vec{r}_i}{\|\vec{r}_i\|^2}$ \quad
  $\vec{F}_{\text{ali}} \propto \overline{\vec{v}}_{\text{voisins}}$ \quad
  $\vec{F}_{\text{coh}} \propto \overline{\vec{p}}_{\text{voisins}} - \vec{p}_{\text{self}}$
};

\end{tikzpicture}

\caption{Forces de séparation, alignement et cohésion appliquées à un boid.}
\label{fig:boids-forces}
\end{figure}

Toutes les forces partagent le même \emph{pattern} de calcul,
ce qui se reflète directement dans le code de
\code{src/entity.c} (chapitre~\ref{chap:implementation})~:

\begin{enumerate}
  \item Parcourir les voisins compatibles dans un rayon donné.
  \item Accumuler un vecteur \emph{désiré} selon la règle.
  \item Calculer un vecteur de \emph{steering} = désiré $-$
    vitesse courante, borné par \code{max\_force}.
  \item Pondérer par le poids de la force et accumuler dans
    l'accélération totale.
\end{enumerate}

L'accélération totale est ensuite bornée à
$4 \times \code{max\_force}$ pour absorber les pics, intégrée
dans la vitesse, qui est elle-même bornée à \code{max\_speed}.

\section{Choix de conception remarquables}

\paragraph{Position centre-based.} La position d'une entité est
le centre de son sprite, pas son coin haut-gauche. Le coin de
dessin est dérivé au moment du rendu. Cela simplifie tous les
calculs de distance et de collision.

\paragraph{Bouche dérivée du cap.} Le point de bouche est
calculé par
\code{pos + sprite\_width × mouth\_forward\_ratio × dir}, où
\code{dir} est le vecteur unitaire de la vitesse. Cela rend
l'alimentation visuellement crédible (le poisson doit \emph{faire
face} à sa proie pour la manger), sans coût de calcul
significatif.

\paragraph{Indices d'espèces dynamiques.} À l'origine, les
espèces étaient identifiées par une \code{enum SpeciesKind}
définie au moment de la compilation. Cette approche figeait la
liste des espèces. Le passage à un index entier dynamique
(0..\code{species\_count}-1) résolu par
\func{species\_find\_by\_id} a permis de découpler la liste des
espèces du code C~: ajouter une espèce ne demande plus de
toucher à \file{species.h}.

\paragraph{Bitmask 32 bits pour \code{prey\_mask}.} Avec
\code{MAX\_SPECIES = 32}, un \code{int} suffit pour décrire
l'ensemble des proies d'un prédateur. Le test
``cette entité $j$ est-elle une proie de l'entité $i$~?'' se
réduit alors à
\code{prey\_mask\_i \& (1 << species\_j)}, donc à une opération
en temps constant. C'est essentiel car ce test est exécuté à
chaque couple $(i, j)$ d'entités, soit en $O(N^2)$ par tick.

\paragraph{Vendoring de \code{tomlc99}.} La bibliothèque TOML
n'est pas une dépendance système~: elle est copiée dans
\file{vendor/tomlc99/} et compilée avec le projet. Cela
garantit la reproductibilité du \emph{build} et l'absence de
téléchargement réseau au moment de la compilation.
